\section{AHSME 1987}
        \begin{enumerate}
        	\item (\textit{AHSME 1987}) $(1+x^2)(1-x^3)$ equals


 $\text{(A)}\ 1 - x^5\qquad \text{(B)}\ 1 - x^6\qquad \text{(C)}\ 1+ x^2 -x^3\qquad \\  \text{(D)}\ 1+x^2-x^3-x^5\qquad \text{(E)}\ 1+x^2-x^3-x^6$


 
 \textbf{Answer: }D



\textbf{Solution}

We multiply: $(1+x^2)(1-x^3) = 1 - x^3 + x^2 - x^5$. Thus the answer is $\boxed{D}$.-slackroadia

 


	\item (\textit{AHSME 1987}) A triangular corner with side lengths $DB=EB=1$ is cut from equilateral triangle ABC of side length $3$. 
The perimeter of the remaining quadrilateral is


 
                    \begin{figure}[H]
                        centering
                        \includegraphics[width=0.5\textwidth]{Problem/AHSME_I_1987_2_Problem_Figure_1.png}
                    \end{figure}
                    


 $\text{(A)} \ 6 \qquad  \text{(B)} \ 6\frac{1}{2} \qquad  \text{(C)} \ 7 \qquad  \text{(D)} \ 7\frac{1}{2} \qquad  \text{(E)} \ 8$


 
 \textbf{Answer: }E



\textbf{Solution}

$\triangle DBE$ is similar to $\triangle ABC$ by AA, so $\overline{DE}$ = 1 by similarity, and $\overline{CE} = \overline{AD} = 2$, by subtraction. Thus the perimeter is $3+2+2+1 = 8$, or $\boxed{E}$. -slackroadia

 


	\item (\textit{AHSME 1987}) How many primes less than $100$ have $7$ as the ones digit? (Assume the usual base ten representation) 


 $\text{(A)} \ 4 \qquad  \text{(B)} \ 5 \qquad  \text{(C)} \ 6 \qquad  \text{(D)} \ 7 \qquad  \text{(E)} \ 8$


 
 \textbf{Answer: }C



\textbf{Solution}

List out all numbers that have 7 as the ones digit less than 100: ${7, 17, 27, 37, 47, 57, 67, 77, 87, 97}$. Only $7, 17,37, 47,67,$ and $97$ are prime. Thus, it is $\boxed{C}$. -slackroadia

 


	\item (\textit{AHSME 1987}) $\frac{2^1+2^0+2^{-1}}{2^{-2}+2^{-3}+2^{-4}}$ equals 


 $\text{(A)} \ 6 \qquad  \text{(B)} \ 8 \qquad  \text{(C)} \ \frac{31}{2} \qquad  \text{(D)} \ 24 \qquad  \text{(E)} \ 512$


 
 \textbf{Answer: }B



\textbf{Solution}

We can factorise to give $\frac{2^1 + 2^0 + 2^{-1}}{2^{-3}(2^1 + 2^0 + 2^{-1})} = \frac{1}{2^{-3}} = 8$, which is $\boxed{\text{B}}$.

 


	\item (\textit{AHSME 1987}) A student recorded the exact percentage frequency distribution for a set of measurements, as shown below. 
However, the student neglected to indicate $N$, the total number of measurements. What is the smallest possible value of $N$?


 
 \[\begin{tabular}{c c}\text{measured value}&\text{percent frequency}\\ \hline 0 & 12.5\\ 1 & 0\\ 2 & 50\\ 3 & 25\\ 4 & 12.5\\ \hline\ & 100\\ \end{tabular}\]


 $\textbf{(A)}\ 5 \qquad \textbf{(B)}\ 8 \qquad \textbf{(C)}\ 16 \qquad \textbf{(D)}\ 25 \qquad \textbf{(E)}\ 50$


 
 \textbf{Answer: }B



\textbf{Solution}

Note that $12.5\% = \frac{1}{8}$, $25\% = \frac{1}{4}$, and $50\% = \frac{1}{2}$. Thus, since the frequencies must be integers, $N$ must be divisible by $2$, $4$, and $8$ (so that $\frac{N}{8}$ etc. are integers), or in other words, $N$ is divisible by $8$. Thus the smallest possible value of $N$ is the smallest positive multiple of $8$, which is $8$ itself, or $\boxed{\text{B}}$.

 


	\item (\textit{AHSME 1987}) In the $\triangle ABC$ shown, $D$ is some interior point, and $x, y, z, w$ are the measures of angles
in degrees. Solve for $x$ in terms of $y, z$ and $w$.


 
                    \begin{figure}[H]
                        centering
                        \includegraphics[width=0.5\textwidth]{Problem/AHSME_I_1987_6_Problem_Figure_1.png}
                    \end{figure}
                    


 $\textbf{(A)}\ w-y-z \qquad \textbf{(B)}\ w-2y-2z \qquad \textbf{(C)}\ 180-w-y-z \qquad \\ \textbf{(D)}\ 2w-y-z\qquad \textbf{(E)}\ 180-w+y+z$


 
 \textbf{Answer: }A



\textbf{Solution}

By angles in a quadrilateral, $x = 360^{\circ} - \text{reflex } \angle ADB - y - z$, and by angles at a point, $\text{reflex } \angle ADB = 360^{\circ} - w$, so our expression becomes $360^{\circ} - (360^{\circ} - w) - y - z = w - y - z$, which is $\boxed{\text{A}}$.

 


	\item (\textit{AHSME 1987}) If $a-1=b+2=c-3=d+4$, which of the four quantities $a,b,c,d$ is the largest?


 $\textbf{(A)}\ a \qquad \textbf{(B)}\ b \qquad \textbf{(C)}\ c \qquad \textbf{(D)}\ d \qquad \textbf{(E)}\ \text{no one is always largest}$


 
 \textbf{Answer: }C



\textbf{Solution}

We have $a = a$, $b = a - 3$, $c = a + 2$, and $d = a - 5$, so now as $2 > 0 > -3 > -5$, $c$ is always largest. Hence the answer is $\boxed{\text{C}}$.

 


	\item (\textit{AHSME 1987}) In the figure the sum of the distances $AD$ and $BD$ is


 
                    \begin{figure}[H]
                        centering
                        \includegraphics[width=0.5\textwidth]{Problem/AHSME_I_1987_8_Problem_Figure_1.png}
                    \end{figure}
                    


 $\textbf{(A)}\ \text{between 10 and 11} \qquad \textbf{(B)}\ 12 \qquad \textbf{(C)}\ \text{between 15 and 16}\qquad\\ \textbf{(D)}\ \text{between 16 and 17}\qquad \textbf{(E)}\ 17$


 
 \textbf{Answer: }C



\textbf{Solution}

Using Pythagoras' Theorem, $BD = \sqrt{3^2 + 4^2} = 5$, and $AD = \sqrt{(13 - 3)^2 + 4^2} = \sqrt{116}$. Thus the sum is $5 + \sqrt{116}$, and as $100 < 116 < 121$, $10 < \sqrt{116} < 11$, so that the sum is between $5+10$ and $5+11$, or $15$ and $16$, which is answer $\boxed{\text{C}}$.

 


	\item (\textit{AHSME 1987}) The first four terms of an arithmetic sequence are $a, x, b, 2x$. The ratio of $a$ to $b$ is


 $\textbf{(A)}\ \frac{1}{4} \qquad \textbf{(B)}\ \frac{1}{3} \qquad \textbf{(C)}\ \frac{1}{2} \qquad \textbf{(D)}\ \frac{2}{3}\qquad \textbf{(E)}\ 2$


 
 \textbf{Answer: }B



\textbf{Solution}

To get from the second term to the fourth term, we add $2$ lots of the common difference, so if the common difference is $d$, we have $2d = 2x -x = x \implies d = \frac{x}{2}$. Thus the sequence is $\frac{x}{2}, x, \frac{3x}{2}, 2x$, so the ratio is $\frac{\frac{x}{2}}{\frac{3x}{2}} = \frac{1}{3}$, which is answer $\boxed{\text{B}}$.

 


	\item (\textit{AHSME 1987}) How many ordered triples $(a, b, c)$ of non-zero real numbers have the property that each number is the product of the other two?


 $\textbf{(A)}\ 1 \qquad \textbf{(B)}\ 2 \qquad \textbf{(C)}\ 3 \qquad \textbf{(D)}\ 4 \qquad \textbf{(E)}\ 5$


 
 \textbf{Answer: }D



\textbf{Solution}

We have $ab = c$, $bc = a$, and $ca = b$, so multiplying these three equations together gives $a^{2}b^{2}c^{2} = abc \implies abc(abc-1)=0$, and as $a$, $b$, and $c$ are all non-zero, we cannot have $abc = 0$, so we must have $abc = 1$. Now substituting $bc = a$ gives $a(bc) = 1 \implies a^2 = 1 \implies a = \pm 1$. If $a = 1$, then the system becomes $b = c, bc = 1, c = b$, so either $b = c = 1$ or $b = c = -1$, giving $2$ solutions. If $a = -1$, the system becomes $-b = c, bc = -1, -c = b$, so $-b = c = 1$ or $b = -c = 1$, giving another $2$ solutions. Thus the total number of solutions is $2 + 2 = 4$, which is answer $\boxed{\text{D}}$.

 


	\item (\textit{AHSME 1987}) Let $c$ be a constant. The simultaneous equations

 \begin{align*}x-y = &\ 2 \\cx+y = &\ 3 \\\end{align*}

have a solution $(x, y)$ inside Quadrant I if and only if


 $\textbf{(A)}\ c=-1 \qquad \textbf{(B)}\ c>-1 \qquad \textbf{(C)}\ c<\frac{3}{2} \qquad \textbf{(D)}\ 0<c<\frac{3}{2}\\ \qquad \textbf{(E)}\ -1<c<\frac{3}{2}$


 
 \textbf{Answer: }E



\textbf{Solution}

We can easily solve the equations algebraically to deduce $x = \frac{5}{c+1}$ and $y = \frac{3-2c}{c+1}$. Thus we firstly need $x > 0 \implies c + 1 > 0 \implies c > -1$. Now $y > 0$ implies $\frac{3-2c}{c+1} > 0$, and since we now know that $c+1$ must be $>0$, the inequality simply becomes $3-2c > 0 \implies 3 > 2c \implies c < \frac{3}{2}$. Thus we combine the inequalities $c > -1$ and $c < \frac{3}{2}$ to get $-1 < c < \frac{3}{2}$, which is answer $\boxed{\text{E}}$.

 


	\item (\textit{AHSME 1987}) In an office, at various times during the day the boss gives the secretary a letter to type, each time putting the letter on top of the pile in the secretary's in-box. When there is time, the secretary takes the top letter off the pile and types it. If there are five letters in all, and the boss delivers them in the order $1\ 2\ 3\ 4\ 5$, which of the following could not be the order in which the secretary types them?


 $\textbf{(A)}\ 1\ 2\ 3\ 4\ 5 \qquad \textbf{(B)}\ 2\ 4\ 3\ 5\ 1 \qquad \textbf{(C)}\ 3\ 2\ 4\ 1\ 5 \qquad \textbf{(D)}\ 4\ 5\ 2\ 3\ 1 \qquad \textbf{(E)}\ 5\ 4\ 3\ 2\ 1 \qquad$


 
 \textbf{Answer: }D



\textbf{Solution}

We can show that $A$, $B$, $C$, and $E$ are possible.

 $A$: the boss delivers $1$, then the secretary does $1$, then the boss delivers $2$, then the secretary does $2$, etc. up to $5$.

 $B$: the boss delivers $1$, then delivers $2$, then the secretary does $2$, then the boss delivers $3$, then delivers $4$, then the secretary does $4$, then the secretary does $3$, then the boss delivers $5$, then the secretary does $5$, then the secretary does $1$.

 $C$: the boss delivers $1$, then delivers $2$, then delivers $3$, then the secretary does $3$, then does $2$, then the boss delivers $4$, then the secretary does $4$, then does $1$, then the boss delivers $5$, then the secretary does $5$.

 $E$: the boss delivers $1$, then $2$, then $3$, then $4$, then $5$, then the secretary does $5$, does $4$, etc. down to $1$.

 Hence the answer must be $\boxed{\text{D}}$. To give a complete proof, we need to show $D$ is impossible. In order for the secretary to start with $4$, the boss must deliver $1$, $2$, $3$, and $4$, then the secretary does $4$. To get $5$ next, the boss next delivers $5$, then the secretary does $5$, but now we can't get $2$, $3$, $1$: the remaining letters must start with $3$, which is on the top of the pile. Hence $D$ is impossible, as required.

 


	\item (\textit{AHSME 1987}) A long piece of paper $5$ cm wide is made into a roll for cash registers by wrapping it $600$ times around a cardboard tube of diameter $2$ cm, 
forming a roll $10$ cm in diameter. Approximate the length of the paper in meters. 
(Pretend the paper forms $600$ concentric circles with diameters evenly spaced from $2$ cm to $10$ cm.)


 $\textbf{(A)}\ 36\pi \qquad \textbf{(B)}\ 45\pi \qquad \textbf{(C)}\ 60\pi \qquad \textbf{(D)}\ 72\pi \qquad \textbf{(E)}\ 90\pi$


 
 \textbf{Answer: }A



\textbf{Solution 1}

Notice (by imagining unfolding the roll), that the length of the paper is equal to the sum of the circumferences of the concentric circles, which is $\pi$ times the sum of the diameters. Now the, the diameters form an arithmetic series with first term $2$, last term $10$, and $600$ terms in total, so using the formula $\frac{1}{2}n(a+l)$, the sum is $300 \times 12 = 3600$, so the length is $3600\pi$ centimetres, or $36\pi$ metres, which is answer $\boxed{\text{A}}$.

 



\textbf{Solution 2}

When the tape is unrolled, its cross-sectional area must remain constant.  The tape's cross-sectional area must be $\pi\cdot 5^2 - \pi \cdot 1^2 = 24\pi$ square centimeters.  Since the roll is $5 - 1 = 4$ cm thick and has $600$ coils, the tape must be $\frac{4}{600} = \frac{1}{150}$ cm thick.  The unrolled tape forms a rectangular prism, so its length must be $\frac{24\pi}{\frac{1}{150}} = 3600\pi$ cm, or $\boxed{(\mathbf{A})\ 36\pi}$ m.

 -j314andrews

 


	\item (\textit{AHSME 1987}) $ABCD$ is a square and $M$ and $N$ are the midpoints of $BC$ and $CD$ respectively. Then $\sin \theta=$


 
                    \begin{figure}[H]
                        centering
                        \includegraphics[width=0.5\textwidth]{Problem/AHSME_I_1987_14_Problem_Figure_1.png}
                    \end{figure}
                    


 $\textbf{(A)}\ \frac{\sqrt{5}}{5} \qquad \textbf{(B)}\ \frac{3}{5} \qquad \textbf{(C)}\ \frac{\sqrt{10}}{5} \qquad \textbf{(D)}\ \frac{4}{5}\qquad \textbf{(E)}\ \text{none of these}$


 
 \textbf{Answer: }B



\textbf{Solution 1}

Use the Sine Area Formula. We can isolate the triangle for which the angle $\theta$ is contained in. WLOG, denote the side length of a triangle as $2$. Our midpoints are then $1$. Subtract the areas of the triangles that don't include the area of our desired triangle: $4-1-1-\frac{1}{2} = \frac{3}{2}.$ The Sine Area Formula tells us $\frac{1}{2}(\sqrt{5})^2\sin\theta=\frac{3}{2}.$ Solving this equation, we get $\sin\theta=\textbf{(B)}\ \frac{3}{5} \qquad$

 Solution: Everyoneintexas

 



\textbf{Solution 2}

Let $\alpha = \angle MAB = \angle DAN$, and let $s$ be the side length of $ABCD$.  Then $BM = \frac{s}{2}$ and $AM = \frac{s\sqrt{5}}{2}$, so $\sin \alpha = \frac{1}{\sqrt{5}}$.  Therefore, $\sin \theta = \sin \left(\frac{\pi}{2} - 2\alpha \right) = \cos 2\alpha = 1 - 2 \sin^{2}\alpha = 1 - 2\cdot \frac{1}{5} = \boxed{\textbf{(B)}\ \frac{3}{5}}$.

 -j314andrews

 


	\item (\textit{AHSME 1987}) If $(x, y)$ is a solution to the system
$xy=6$ and $x^2y+xy^2+x+y=63$,
find $x^2+y^2$.


 $\textbf{(A)}\ 13 \qquad \textbf{(B)}\ \frac{1173}{32} \qquad \textbf{(C)}\ 55 \qquad \textbf{(D)}\ 69 \qquad \textbf{(E)}\ 81$


 
 \textbf{Answer: }D



\textbf{Solution}

First note that $x^2y+xy^2+x+y= (xy+1)(x+y)$. Substituting $6$ for $xy$ gives $7(x+y)= 63$, giving a result of $x+y=9$. Squaring this equation and subtracting by $12$, gives us $x^2+y^2= \boxed{69}$

 


	\item (\textit{AHSME 1987}) A cryptographer devises the following method for encoding positive integers. First, the integer is expressed in base $5$. 
Second, a 1-to-1 correspondence is established between the digits that appear in the expressions in base $5$ and the elements of the set 
$\{V, W, X, Y, Z\}$. Using this correspondence, the cryptographer finds that three consecutive integers in increasing 
order are coded as $VYZ, VYX, VVW$, respectively. What is the base-$10$ expression for the integer coded as $XYZ$?


 $\textbf{(A)}\ 48 \qquad \textbf{(B)}\ 71 \qquad \textbf{(C)}\ 82 \qquad \textbf{(D)}\ 108 \qquad \textbf{(E)}\ 113$


 
 \textbf{Answer: }D



\textbf{Solution}

Since $VYX + 1 = VVW$, i.e. adding $1$ causes the "fives" digit to change, we must have $X = 4$ and $W = 0$. Now since $VYZ + 1 = VYX$, we have $X = Z + 1 \implies Z = 4 - 1 = 3$. Finally, note that in $VYX + 1 = VVW$, adding $1$ will cause the "fives" digit to change by $1$ if it changes at all, so $V = Y + 1$, and thus since $1$ and $2$ are the only digits left (we already know which letters are assigned to $0$, $3$, and $4$), we must have $V = 2$ and $Y = 1$. Thus $XYZ = 413_{5} = 4 \cdot 5^{2} + 1 \cdot 5 + 3 = 100 + 5 + 3 = 108$, which is answer $\boxed{\text{D}}$.

 


	\item (\textit{AHSME 1987}) In a mathematics competition, the sum of the scores of Bill and Dick equalled the sum of the scores of Ann and Carol. 
If the scores of Bill and Carol had been interchanged, then the sum of the scores of Ann and Carol would have exceeded 
the sum of the scores of the other two. Also, Dick's score exceeded the sum of the scores of Bill and Carol. 
Determine the order in which the four contestants finished, from highest to lowest. Assume all scores were nonnegative.


 $\textbf{(A)}\ \text{Dick, Ann, Carol, Bill} \qquad \textbf{(B)}\ \text{Dick, Ann, Bill, Carol} \qquad \textbf{(C)}\ \text{Dick, Carol, Bill, Ann}\\ \qquad \textbf{(D)}\ \text{Ann, Dick, Carol, Bill}\qquad \textbf{(E)}\ \text{Ann, Dick, Bill, Carol}$


 
 \textbf{Answer: }E



\textbf{Solution}

Let the scores be $A$, $B$, $C$, and $D$. Then we have $A + C = B + D$, $A + B > C + D$, and $D > B + C$. Call these equations $(1)$, $(2)$, and $(3)$ respectively. Then adding $(1)$ and $(2)$ gives $2A + B + C > B + C + 2D \implies 2A > 2D \implies A > D$. Now subtracting $(1)$ from $(2)$ gives $A + B - A - C > C + D - B - D \implies B - C > C - B \implies 2B > 2C \implies B > C$. Finally from $(3)$, since  $D > B + C$ and $C$ is non-negative, we must have $D > B$, so putting our results together we get $A > D > B > C$, which is answer $\boxed{\text{E}}$.

 


	\item (\textit{AHSME 1987}) It takes $A$ algebra books (all the same thickness) and $H$ geometry books (all the same thickness, 
which is greater than that of an algebra book) to completely fill a certain shelf. 
Also, $S$ of the algebra books and $M$ of the geometry books would fill the same shelf. 
Finally, $E$ of the algebra books alone would fill this shelf. Given that $A, H, S, M, E$ are distinct positive integers, 
it follows that $E$ is


 $\textbf{(A)}\ \frac{AM+SH}{M+H} \qquad \textbf{(B)}\ \frac{AM^2+SH^2}{M^2+H^2} \qquad \textbf{(C)}\ \frac{AH-SM}{M-H}\qquad \textbf{(D)}\ \frac{AM-SH}{M-H}\qquad \textbf{(E)}\ \frac{AM^2-SH^2}{M^2-H^2}$


 
 \textbf{Answer: }D



\textbf{Solution}

Let $x$ and $y$ be the thicknesses of an algebra book and geometry book, respectively, and let $z$ be the length of the shelf. Then from the given information,
 \[Ax + Hy = z\]
 \[Sx + My = z\]
 \[Ex = z.\] From the third equation, $x = z/E$. Substituting into the first two equations, we get
 \[\frac{A}{E} z + Hy = z,\] 
 \[\frac{S}{E} z + My = z.\] 

 From the first equation,
 \[Hy = z - \frac{A}{E} z = \frac{E - A}{E} z,\]so
 \[\frac{y}{z} = \frac{E - A}{EH}.\]From the second equation,
 \[My = z - \frac{S}{E} z = \frac{E - S}{E} z,\]so
 \[\frac{y}{z} = \frac{E - S}{ME}.\]Hence,
 \[\frac{E - A}{EH} = \frac{E - S}{ME}.\]Multiplying both sides by $HME$, we get $ME - AM = HE - HS$. Then $(M - H)E = AM - HS$, so
 \[E = \boxed{\frac{AM - HS}{M - H}}.\]The answer is (D).

 



\textbf{Solution 2 (a little faster)}

Let $a$ and $g$ denote the widths of an algebra and geometry textbook, respectively. As in Solution 1, observe that the width of the shelf is equal to the following: 

 
 \[Ea=Aa+Hg=Sa+Mg.\]

 The rest is simply algebra: 

 \begin{align*}(E-A)a&=Hg \\g&=\dfrac{(E-A)a}H \\Ea&=Sa+M\left(\dfrac{(E-A)a}H\right) \\&=Sa+\dfrac{M(E-A)a}H \\E&=S+\dfrac{M(E-A)}H \\&=\dfrac{SH+ME-MA}H \\EH&=SH+ME-MA \\E(H-M)&=SH-MA \\E&=\dfrac{SH-MA}{H-M} \\&=\dfrac{AM-SH}{M-H}. \\\end{align*}


 This is answer choice $\boxed{\textbf{(D)}}$. QED. $\Box$

 


	\item (\textit{AHSME 1987}) Which of the following is closest to $\sqrt{65}-\sqrt{63}$?


 $\textbf{(A)}\ .12 \qquad \textbf{(B)}\ .13 \qquad \textbf{(C)}\ .14 \qquad \textbf{(D)}\ .15 \qquad \textbf{(E)}\ .16$


 
 \textbf{Answer: }B



\textbf{Solution}

We have $\sqrt{65} > 8 > 7.5$. Also $7.5^2 = (7 + 0.5)^2 = 7^2 + 2 \cdot 7 \cdot 0.5 + 0.5^2 = 49 + 7 + 0.25 = 56.25 < 63$, so $\sqrt{63} > 7.5$. Thus $\sqrt{65} + \sqrt{63} > 7.5 + 7.5 = 15$. Now notice that $\sqrt{65} - \sqrt{63} = \frac{(\sqrt{65} - \sqrt{63})(\sqrt{65} + \sqrt{63})}{\sqrt{65} + \sqrt{63}} = \frac{2}{\sqrt{65} + \sqrt{63}}$, so $\sqrt{65} - \sqrt{63} < \frac{2}{15} = 0.1333333...$, so the answer must be $A$ or $B$. If we can show \sqrt{65} + \sqrt{63} < 16, then we will have $\frac{2}{\sqrt{65} + \sqrt{63}} > 2/16 = 0.125$ and hence the answer must be $B$. 

 
The proof of $\sqrt{65} + \sqrt{63} < 16$ is as follows:

 $\sqrt{65} + \sqrt{63} < 16 \iff 65 + 2\sqrt{65 \cdot 63} + 63 < 256 \iff 128 + 2\sqrt{4095} < 256 \iff \sqrt{4095} < 64 \iff 4095 < 4096$ which is true. 

 
($\sqrt{65} - \sqrt{63}$ is approximately equal to $0.125003815$)

 


	\item (\textit{AHSME 1987}) Evaluate
$\log_{10}(\tan 1^{\circ})+\log_{10}(\tan 2^{\circ})+\log_{10}(\tan 3^{\circ})+\cdots+\log_{10}(\tan 88^{\circ})+\log_{10}(\tan 89^{\circ}).$


 $\textbf{(A)}\ 0 \qquad \textbf{(B)}\ \frac{1}{2}\log_{10}(\frac{\sqrt{3}}{2}) \qquad \textbf{(C)}\ \frac{1}{2}\log_{10}2\qquad \textbf{(D)}\ 1\qquad \textbf{(E)}\ \text{none of these}$


 
 \textbf{Answer: }A



\textbf{Solution}

Because $\tan x \tan (90^\circ - x) = \tan x \cot x = 1$, $\tan 45^\circ = 1$, and $\log a + \log b = \log {ab}$, the answer is $\log_{10} {\tan 1^\circ \tan 2^\circ \dots \tan 89^\circ} = \log_{10} 1 = 0.$ $\boxed{\textbf{(A)}}.$

 



\textbf{Solution 2}

We have $\log{(\frac{\sin 1^\circ}{\cos 1^\circ} \times \frac{\sin 2^\circ}{\cos 2^\circ} \times ... \times \frac{\sin 89^\circ}{\cos 89^\circ})}$. However, $\sin{(x)} = \cos{(90^\circ - x)}$. Thus each pair of $\sin, \cos$ (for example, $\sin{1^\circ}, \cos{89^\circ}$) multiplies to $1$. Hence we have $\log{1} = 0$.

 


	\item (\textit{AHSME 1987}) There are two natural ways to inscribe a square in a given isosceles right triangle. 
If it is done as in Figure 1 below, then one finds that the area of the square is $441 \text{cm}^2$. 
What is the area (in $\text{cm}^2$) of the square inscribed in the same $\triangle ABC$ as shown in Figure 2 below?


 
                    \begin{figure}[H]
                        centering
                        \includegraphics[width=0.5\textwidth]{Problem/AHSME_I_1987_21_Problem_Figure_1.png}
                    \end{figure}
                    


 $\textbf{(A)}\ 378 \qquad \textbf{(B)}\ 392 \qquad \textbf{(C)}\ 400 \qquad \textbf{(D)}\ 441 \qquad \textbf{(E)}\ 484$


 
 \textbf{Answer: }B



\textbf{Solution}

We are given that the area of the inscribed square is $441$, so the side length of that square is $21$. Since the square divides the $45-45-90$ larger triangle into 2 smaller congruent $45-45-90$, then the legs of the larger isosceles right triangle ($BC$ and $AB$) are equal to $42$. 
                    \begin{figure}[H]
                        centering
                        \includegraphics[width=0.5\textwidth]{Solution/AHSME_I_1987_21_Solution_Figure_1.png}
                    \end{figure}
                    

 We now have that $3S=42\sqrt{2}$, so $S=14\sqrt{2}$. But we want the area of the square which is $S^2=(14\sqrt{2})^2= \boxed{\mathrm{(B)}\ 392}$

 


	\item (\textit{AHSME 1987}) A ball was floating in a lake when the lake froze. The ball was removed (without breaking the ice), 
leaving a hole $24$ cm across as the top and $8$ cm deep. What was the radius of the ball (in centimeters)?


 $\textbf{(A)}\ 8 \qquad \textbf{(B)}\ 12 \qquad \textbf{(C)}\ 13 \qquad \textbf{(D)}\ 8\sqrt{3} \qquad \textbf{(E)}\ 6\sqrt{6}$


 
 \textbf{Answer: }C



\textbf{Solution}

Consider a cross-section of this problem in which a circle lies with its center somewhere above a line. A line segment of $8$ cm can be drawn from the line to the bottom of the ball. Denote the distance between the center of the circle and the line as $x$. We can construct a right triangle by dragging the center of the circle to the intersection of the circle and the line. We then have the equation $x^2+(12)^2=(x+8)^2$, $x^2+144=x^2+16x+64$. Solving, the answer is ${\textbf{(C)}\ 13}$.

 


	\item (\textit{AHSME 1987}) If $p$ is a prime and both roots of $x^2+px-444p=0$ are integers, then


 $\textbf{(A)}\ 1<p\le 11 \qquad \textbf{(B)}\ 11<p \le 21 \qquad \textbf{(C)}\ 21< p \le 31 \\ \qquad \textbf{(D)}\ 31< p\le 41\qquad \textbf{(E)}\ 41< p\le 51$


 
 \textbf{Answer: }D



\textbf{Solution}

For integer roots, we need the discriminant, which is $p^2 - 4 \cdot 1 \cdot (-444p) = p^2 + 1776p = p(p+1776)$, to be a perfect square. Now, this means that $p$ must divide $p + 1776$, as if it did not, there would be a lone prime factor of $p$, and so this expression could not possibly be a perfect square. Thus $p$ divides $p + 1776$, which implies $p$ divides $1776 = 2^{4} \cdot 3 \cdot 37$, so we must have $p = 2$, $3$, or $37$. It is easy to verify that neither $p = 2$ nor $p = 3$ make $p(p+1776)$ a perfect square, but $p = 37$ does, so the answer is $31 < p \le 41$, which is $\boxed{\text{D}}$.

 


	\item (\textit{AHSME 1987}) How many polynomial functions $f$ of degree $\ge 1$ satisfy
$f(x^2)=[f(x)]^2=f(f(x))$  ?


 $\textbf{(A)}\ 0 \qquad \textbf{(B)}\ 1 \qquad \textbf{(C)}\ 2 \qquad \textbf{(D)}\ \text{finitely many but more than 2}\\ \qquad \textbf{(E)}\ \infty$


 
 \textbf{Answer: }B



\textbf{Solution}

Let $f(x) = \sum_{k=0}^{n} a_{k} x^{k}$ be a polynomial satisfying the condition, so substituting it in, we find that the highest powers in each of the three expressions are, respectively, $a_{n}x^{2n}$, $a_{n}^{2}x^{2n}$, and $a_{n}^{n+1}x^{n^{2}}$. If polynomials are identically equal, each term must be equal, so we get $2n = n^2$ and $a_{n} = a_{n}^{2} = a_{n}^{n+1}$, so since $n \geq 1$, we must have $n = 2$, and since $a_{n} \neq 0$, we have $a_{n} = 1$. The given condition now becomes $x^4 + bx^2 + c \equiv (x^2 + bx + c)^2$, so we must have $b = 0$, or else the right-hand side would have a cubic term that the left-hand side does not. Thus we get $x^4 + c \equiv (x^2 + c)^2$, so we must have $c = 0$, or else the right-hand side would have an $x^2$ term that the left-hand side does not. Thus the only possibility is $f(x) = x^2$, i.e. there is only $1$ solution, so the answer is $\boxed{\text{B}}$.

 


	\item (\textit{AHSME 1987}) $ABC$ is a triangle: $A=(0,0), B=(36,15)$ and both the coordinates of $C$ are integers. What is the minimum area $\triangle ABC$ can have?


 $\textbf{(A)}\ \frac{1}{2} \qquad \textbf{(B)}\ 1 \qquad \textbf{(C)}\ \frac{3}{2} \qquad \textbf{(D)}\ \frac{13}{2}\qquad \textbf{(E)}\ \text{there is no minimum}$


 
 \textbf{Answer: }C



\textbf{Solution}

Let $C$ have coordinates $(p, q)$. Then by the Shoelace Formula, the area of $\triangle ABC$ is $\frac{3}{2} \lvert {12q-5p} \rvert$. Since $p$ and $q$ are integers, $\lvert {12q-5p} \rvert$ is a positive integer, and by Bezout's Lemma, it can equal $1$ (e.g. with $q = 2, p = 5$), so the minimum area is $\frac{3}{2} \times 1 = \frac{3}{2}$, which is answer $\boxed{C}$.

 


 



\textbf{Solution 2}

Let $C$ have coordinates $(p, q)$. The length of $AB$ is $\sqrt{(36-0)^2+(15-0)^2}=39.$ The equation of line $AB$ is $5x-12y=0.$ Then, by the point-line distance formula the length of the altitude from $C$ to $AB$ is $\frac{|5p-12q|}{\sqrt{5^2+12^2}}=\frac{|5p-12q|}{13}.$ Thus the area of the triangle is $39 \cdot \frac{|5p-12q|}{13}=\frac{3}{2} \lvert {12q-5p} \rvert.$ Then proceed as in solution $1$ from here.

 -AndrewZhong2012

 


	\item (\textit{AHSME 1987}) The amount $2.5$ is split into two nonnegative real numbers uniformly at random, 
for instance, into $2.143$ and $.357$, or into $\sqrt{3}$ and $2.5-\sqrt{3}$. 
Then each number is rounded to its nearest integer, for instance, $2$ and $0$ in the first case above, 
$2$ and $1$ in the second. What is the probability that the two integers sum to $3$?


 $\textbf{(A)}\ \frac{1}{4} \qquad \textbf{(B)}\ \frac{2}{5} \qquad \textbf{(C)}\ \frac{1}{2} \qquad \textbf{(D)}\ \frac{3}{5}\qquad \textbf{(E)}\ \frac{3}{4}$


 
 \textbf{Answer: }B



\textbf{Solution}

The two parts may round to $0$ and $2$; $1$ and $2$; $1$ and $1$; $2$ and $1$; or $2$ and $0$. By considering the possible ranges for each case, it is easy to see that each case is equally likely (they divide the interval from $0$ to $2.5$, in which one of the parts is found, into five equal ranges of $0$ to $0.5$, $0.5$ to $1$, $1$ to $1.5$, $1.5$ to $2$, and $2$ to $2.5$). As exactly two of the five cases give a sum of $1 + 2 = 3$, the answer is $\frac{2}{5}$, which is answer $\boxed{B}$.

 


	\item (\textit{AHSME 1987}) A cube of cheese $C=\{(x, y, z)| 0 \le x, y, z \le 1\}$ is cut along the planes $x=y, y=z$ and $z=x$. How many pieces are there? 
(No cheese is moved until all three cuts are made.)


 $\textbf{(A)}\ 5 \qquad \textbf{(B)}\ 6 \qquad \textbf{(C)}\ 7 \qquad \textbf{(D)}\ 8 \qquad \textbf{(E)}\ 9$


 
 \textbf{Answer: }B



\textbf{Solution}

The cut $x = y$ separates the cube into points with $x < y$ and points with $x > y$, and analogous results apply for the other cuts. Thus, which piece a particular point is in depends only on the relative sizes of its coordinates $x$, $y$, and $z$ - for example, all points with the ordering $x < y < z$ are in the same piece. Thus, as there are $3! = 6$ possible orderings, there are $6$ pieces, which is answer $\boxed{B}$.

 


	\item (\textit{AHSME 1987}) Let $a, b, c, d$ be real numbers. Suppose that all the roots of $z^4+az^3+bz^2+cz+d=0$ are complex numbers lying on a circle
in the complex plane centered at $0+0i$ and having radius $1$. The sum of the reciprocals of the roots is necessarily


 $\textbf{(A)}\ a \qquad \textbf{(B)}\ b \qquad \textbf{(C)}\ c \qquad \textbf{(D)}\ -a \qquad \textbf{(E)}\ -b$


 
 \textbf{Answer: }D



\textbf{Solution}

Let's denote the roots of the polynomial as $z_1, z_2, z_3, z_4$. We know that the magnitudes of these 4 roots are 1 as given in the problem statement. Therefore, we have $z_1 \overline{z_1}, z_2 \overline{z_2}, z_3 \overline{z_3}, z_4 \overline{z_4} = 1$. We want to find $\frac{1}{z_1} + \frac{1}{z_2} + \frac{1}{z_3} + \frac{1}{z_4}$ which is $\overline{z_1} + \overline{z_2} + \overline{z_3} + \overline{z_4}$. Remember that these are the roots of polynomials. Whenever complex numbers are the roots of a polynomial with real coefficients, we know they come in complex-conjugate pairs. Therefore, $\overline{z_1} + \overline{z_2} + \overline{z_3} + \overline{z_4} = z_1 + z_2 + z_3 + z_4$. By Vieta's, $z_1 + z_2 + z_3 + z_4 = -a$. Thus, the answer is $\boxed{\textbf{(D) } -a }$. 

 --$\text{lucasxia01}$

 


	\item (\textit{AHSME 1987}) Consider the sequence of numbers defined recursively by $t_1=1$ and for $n>1$ by $t_n=1+t_{(n/2)}$ when $n$ is even 
and by $t_n=\frac{1}{t_{(n-1)}}$ when $n$ is odd. Given that $t_n=\frac{19}{87}$, the sum of the digits of $n$ is


 $\textbf{(A)}\ 15 \qquad \textbf{(B)}\ 17 \qquad \textbf{(C)}\ 19 \qquad \textbf{(D)}\ 21 \qquad \textbf{(E)}\ 23$


 
 \textbf{Answer: }A



\textbf{Solution}

If $n$ is even, then $t_{(n/2)}$ would be negative, which is not possible. Therefore, $n$ is odd. With this function, backwards thinking is the key. If $t_x < 1$, then $x$ is odd, and $t_{(x-1)} = \frac{1}{t_{x}}$. Otherwise, you keep on subtracting 1 and halving x until $t_\frac{x}{2^{n}} < 1$.We can use this logic to go backwards until we reach $t_1 = 1$, like so:

 $t_n=\frac{19}{87}\\\\t_{n-1} = \frac{87}{19}\\\\t_{\frac{n-1}{2}} = \frac{68}{19}\\\\t_{\frac{n-1}{4}} = \frac{49}{19}\\\\t_{\frac{n-1}{8}} = \frac{30}{19}\\\\t_{\frac{n-1}{16}} = \frac{11}{19}\\\\t_{\frac{n-1}{16} - 1} = \frac{19}{11}\\\\t_{\frac{\frac{n-1}{16} - 1}{2}} = \frac{8}{11}\\\\t_{\frac{\frac{n-1}{16} - 1}{2} - 1} = \frac{11}{8}\\\\t_{\frac{\frac{\frac{n-1}{16} - 1}{2} - 1}{2}} = \frac{3}{8}\\\\t_{\frac{\frac{\frac{n-1}{16} - 1}{2} - 1}{2} - 1} = \frac{8}{3}\\\\t_{\frac{\frac{\frac{\frac{n-1}{16} - 1}{2} - 1}{2} - 1}{2}} = \frac{5}{3}\\\\t_{\frac{\frac{\frac{\frac{n-1}{16} - 1}{2} - 1}{2} - 1}{4}} = \frac{2}{3}\\\\t_{\frac{\frac{\frac{\frac{n-1}{16} - 1}{2} - 1}{2} - 1}{4} - 1} = \frac{3}{2}\\\\t_{\frac{\frac{\frac{\frac{\frac{n-1}{16} - 1}{2} - 1}{2} - 1}{4} - 1}{2}} = \frac{1}{2}\\\\t_{\frac{\frac{\frac{\frac{\frac{n-1}{16} - 1}{2} - 1}{2} - 1}{4} - 1}{2} - 1} = 2\\\\t_{\frac{\frac{\frac{\frac{\frac{\frac{n-1}{16} - 1}{2} - 1}{2} - 1}{4} - 1}{2} - 1}{2}} = t_1 = 1 \Rightarrow \frac{\frac{\frac{\frac{\frac{\frac{n-1}{16} - 1}{2} - 1}{2} - 1}{4} - 1}{2} - 1}{2} = 1 \Rightarrow n = 1905$, so the answer is $\boxed{\textbf{(A)} 15}$.

 


	\item (\textit{AHSME 1987}) In the figure, $\triangle ABC$ has $\angle A =45^{\circ}$ and $\angle B =30^{\circ}$. A line $DE$, with $D$ on $AB$ 
and $\angle ADE =60^{\circ}$, divides $\triangle ABC$ into two pieces of equal area. 
(Note: the figure may not be accurate; perhaps $E$ is on $CB$ instead of $AC$.) 
The ratio $\frac{AD}{AB}$ is


 
                    \begin{figure}[H]
                        centering
                        \includegraphics[width=0.5\textwidth]{Problem/AHSME_I_1987_30_Problem_Figure_1.png}
                    \end{figure}
                    


 $\textbf{(A)}\ \frac{1}{\sqrt{2}} \qquad \textbf{(B)}\ \frac{2}{2+\sqrt{2}} \qquad \textbf{(C)}\ \frac{1}{\sqrt{3}} \qquad \textbf{(D)}\ \frac{1}{\sqrt[3]{6}}\qquad \textbf{(E)}\ \frac{1}{\sqrt[4]{12}}$


 
 \textbf{Answer: }E



\textbf{Solution}

First we show that $E$ is on $AC$, as in the given figure, by demonstrating that if $E = C$, then $\triangle ADE$ has more than half the area, so $DE$ is too far to the right. Specifically, assume $E = C$ and drop an altitude from $C$ to $AB$ that meets $AB$ at $F$. Without loss of generality, assume that $CF = 1$, so that $\frac{\text{Area} \ \triangle EAD}{\text{Area} \ \triangle EAB} = \frac{AD}{AB}$ (as they have the same height, and area is $\frac{1}{2}$ times the product of base and height), which is $\frac{1 + \frac{1}{\sqrt{3}}}{1 + \sqrt{3}} = \frac{1}{\sqrt{3}} > \frac{1}{2}$, as required. Thus we must move $DE$ to the left, scaling $\triangle EAD$ by a factor of $k$ such that $\text{Area} \ \triangle EAD = \frac{1}{2} \ \text{Area} \ \triangle CAB \implies \frac{1}{2}k^{2}(1+\frac{1}{\sqrt{3}}) = \frac{1}{4}(1+\sqrt{3}) \implies k^{2} = \frac{\sqrt{3}}{2} \implies k = (\frac{3}{4})^{\frac{1}{4}}$. Thus $\frac{AD}{AB} = \frac{k(1+\frac{1}{\sqrt{3}})}{1+\sqrt{3}} = \frac{k}{\sqrt{3}} = (\frac{3}{4})^{\frac{1}{4}} \times (\frac{1}{9})^{\frac{1}{4}} = (\frac{1}{12})^{\frac{1}{4}} = \frac{1}{\sqrt[4]{12}}$, which is answer $\boxed{E}$.

 



\textbf{Solution 2 (Trigonometry)}

The proof of why point $E$ lies on $AC$ is the same as Solution 1.

 We have:\begin{align}\text{Area} \ \triangle AED &= \text{Area} \ \triangle CDB\\\text{Area} \ \triangle AED &= \text{Area} \ \triangle ABC - \text{Area} \ \triangle AED\\2\cdot\text{Area} \ \triangle AED &= \text{Area} \ \triangle ABC\end{align}Since the area of a triangle can be expressed as $\dfrac 12ab\sin C$, we have\begin{align}2 \cdot \dfrac 12 AE\cdot AD \cdot \sin45^\circ &= \dfrac 12 AC\cdot AB \cdot \sin45^\circ\\2 \cdot AE \cdot AD = AC \cdot AB\end{align}

 Now, by Law of Sines, we have $AE = \dfrac{AD \cdot \sin 60^\circ}{\sin 75^\circ}$ and $AC = \dfrac{AB \cdot \sin 30 ^\circ}{\sin 75^\circ}$. Thus, by substitution we have\begin{align}2 \cdot AD \cdot \dfrac{AD \cdot \sin 60^\circ}{\sin 75^\circ} &= AB \cdot \dfrac{AB \cdot \sin 30 ^\circ}{\sin 75^\circ}\\2 \cdot AD^2 \cdot \dfrac{\sqrt{3}}2 &= AB^2 \cdot \dfrac 12\\\dfrac{AD^2}{AB^2} &= \dfrac{1}{2\sqrt{3}}\\\dfrac{AD}{AB} &= \dfrac{1}{\sqrt[4]{12}}\end{align}The answer is therefore $\boxed{E}$.

 


\end{enumerate}